% ============================================================
%  Remerciements — Mémoire de Master BIBDA — CHOUFA Zouhair
% ============================================================
\clearpage
\thispagestyle{empty}

\begin{center}
    \vspace*{1cm} 
    {\Huge \scshape \color{titleblue} \textbf{Remerciements}} \\[0.2cm]
    \textcolor{bordergreen}{\rule{1cm}{1.2pt}}
    \vspace{1cm} 
\end{center}
\addcontentsline{toc}{chapter}{Remerciements}
\renewcommand{\baselinestretch}{1.15}\normalsize

Au terme de ce travail, il est adressé une gratitude sincère à toutes
les personnes ayant contribué, directement ou indirectement, à la
réalisation de ce mémoire.

Les remerciements les plus sincères sont exprimés au
\textbf{Professeur A. MADANI}, encadrant de ce travail à la Faculté
des Sciences d'El Jadida, pour la qualité de son encadrement, sa
disponibilité constante et la pertinence de ses conseils tout au long
de ce projet. Sa rigueur scientifique et la confiance accordée dans
la conduite de ce sujet de recherche ont permis de mener cette étude
dans les meilleures conditions.

Les membres du jury sont également remerciés pour l'honneur qu'ils
font en évaluant ce travail. Leur lecture attentive et leurs
remarques constitueront un apport précieux à la poursuite de cette
réflexion scientifique.

Une reconnaissance particulière est adressée au corps professoral et
administratif du \textbf{Département d'Informatique} de la Faculté
des Sciences d'El Jadida, ainsi qu'à \\ l'\textbf{Université Chouaib
Doukkali}, pour la qualité de la formation dispensée durant ce
Master, dont les enseignements théoriques et pratiques ont constitué
le socle de ce travail.

Les camarades de la promotion BIBDA sont remerciés pour l'esprit
d'entraide et les échanges techniques qui ont accompagné ce parcours,
en particulier durant les phases les plus intenses de développement
et de rédaction.

Une pensée est également adressée à la communauté scientifique et
technique open source — notamment les contributeurs des projets
\texttt{Morph-KGC}, \texttt{RDFLib}, \texttt{Splink},
\texttt{Sentence-- Transformers} et \texttt{LangChain} — dont les
outils ont rendu possible l'implémentation du pipeline présenté dans
ce mémoire.

Enfin, les remerciements les plus profonds vont à la famille, pour
son soutien indéfectible, ses encouragements constants et les
sacrifices consentis tout au long de ce parcours académique. Sans ce
soutien, ce mémoire n'aurait pu voir le jour.

\newpage